\section{Cumulative-storage gluing and joint continuation}\label{sec:continuation}

This section gives the algebra behind the continuation criterion in
the exposition~\cite{Exposition}. It follows the inherited
cumulative-storage note~\cite{CumulativePath}. All operators here
act on ordinary unweighted interval Hilbert spaces. The identities
supply tests for a proposed construction. The final subsection reports
a separate internally certified fixed-window anchor, a quantitative
obstruction to a scalar coupling estimate, and a subsequent bounded append. No successful all-depth recursion
is established.

\subsection{The fixed-shift spatial Schur complement}

After translating $I_{L+h}$ to $(0,L+h)$, causality gives
\begin{equation}\label{eq:si-block}
 V=V_{\omega,L+h}=
 \begin{pmatrix}X&0\\Y&Z\end{pmatrix},
 \qquad X=V_{\omega,L},\quad Z=V_{\omega,h}.
\end{equation}
Here $Y$ carries old input into the added output interval.
Suppose $\norm X<1$ and $\norm Z<1$, and set
\[
 E=I-X^*X,\qquad F_{\mathrm{out}}=I-ZZ^*,
 \qquad \mathcal C=F_{\mathrm{out}}^{-1/2}YE^{-1/2}.
\]
These defects are coercive: each is bounded below by a positive
multiple of the identity.

\begin{proposition}[Cumulative gluing criterion]\label{prop:si-gluing}
Under these strict diagonal bounds,
\begin{equation}\label{eq:si-gluing}
 \norm V\leq1
 \quad\Longleftrightarrow\quad
 E-Y^*F_{\mathrm{out}}^{-1}Y\succeq0
 \quad\Longleftrightarrow\quad\norm{\mathcal C}\leq1.
\end{equation}
Replacing both norm inequalities by strict inequalities gives
the strict-contraction criterion.
\end{proposition}
\begin{proof}
The defect is
\[
 I-V^*V=
 \begin{pmatrix}
 E-Y^*Y&-Y^*Z\\
 -Z^*Y&I-Z^*Z
 \end{pmatrix}.
\]
Its lower diagonal block is coercive. Its Schur complement is
\[
 E-Y^*\bigl[I+Z(I-Z^*Z)^{-1}Z^*\bigr]Y
 =E-Y^*F_{\mathrm{out}}^{-1}Y
 =E^{1/2}(I-\mathcal C^*\mathcal C)E^{1/2}.
\]
The middle identity follows by multiplying by $I-ZZ^*$.
Schur elimination proves positivity and, using bounded invertible
triangular elimination factors, equivalence of coercivity.
\end{proof}
The output orientation $I-ZZ^*$ is essential; replacing it with
$I-Z^*Z$ inside the displayed cross term changes the condition.
At non-strict diagonal contraction the inverses need not exist,
so this proof cannot be used without a range-compatible extension.

\subsection{Cayley coordinates and reflection}

For $\omega>0$ the finite-window transfer is Volterra convolution
by a locally integrable kernel. It is quasinilpotent: conjugation
by $e^{-cx}$ gives a convolution norm bounded by
$\int_0^L e^{-cu}|k_\omega(u)|\dd u$, which tends to zero as
$c\to\infty$. Each finite-$c$ conjugation is a bounded similarity,
so the spectral radius is zero. Hence $I+V$ has a bounded inverse
without assuming contraction.

Define
\begin{equation}
 \mathcal Z=(I-V)(I+V)^{-1},\qquad
 P_{\omega,L}=\frac2\omega\Ree\mathcal Z,\qquad
 \Ree A=\tfrac12(A+A^*).
\end{equation}
Multiplying the real part on the two sides gives
\begin{equation}\label{eq:si-congruence}
 I-V^*V=\frac\omega2(I+V^*)P_{\omega,L}(I+V).
\end{equation}
Thus $P_{\omega,L}\succeq0$ is exactly contraction; coercivity
is exactly strict contraction in operator norm.
For the real causal kernel, interval reflection
$J_Lf(x)=f(-x)$ on the centered interval gives
$V^*=J_LVJ_L$. Consequently $\mathcal Z^*=J_L\mathcal ZJ_L$,
and $P_{\omega,L}$ commutes with reflection. The raw input
defect need not commute with it.

For the block split \eqref{eq:si-block}, put
$R_X=(I+X)^{-1}$, $R_Z=(I+Z)^{-1}$. Block inversion yields
\begin{equation}
 P_{\omega,L+h}=\begin{pmatrix}A&B^*\\B&F\end{pmatrix},
 \quad A=P_{\omega,L},\quad F=P_{\omega,h},\quad
 B=-\frac2\omega R_ZYR_X.
\end{equation}
The two useful orientations of the diagonal identities are
\begin{equation}
 A=\frac2\omega R_X^*ER_X,\qquad
 F=\frac2\omega R_ZF_{\mathrm{out}}R_Z^*.
\end{equation}
Under the strict bounds of Proposition~\ref{prop:si-gluing},
\begin{equation}\label{eq:si-coupling}
 \norm{F^{-1/2}BA^{-1/2}}
 =\norm{F_{\mathrm{out}}^{-1/2}YE^{-1/2}}.
\end{equation}
Indeed the left normalized block equals $-\mathcal U\mathcal C\mathcal W$,
where
\[
 \mathcal U=\sqrt{2/\omega}\,F^{-1/2}R_ZF_{\mathrm{out}}^{1/2},
 \qquad
 \mathcal W=\sqrt{2/\omega}\,E^{1/2}R_XA^{-1/2}.
\]
The displayed diagonal identities make these invertible maps
unitary. This proves \eqref{eq:si-coupling} without commuting any
noncommuting factors. The Cayley change preserves the full
coupling norm; it provides coordinates rather than stronger bounds.

\subsection{A step changing both length and shift}

Fix the proposed append length $h$. At a reference positive shift
$\alpha$, suppose $A_\alpha=P_{\alpha,L}$ and
$F_\alpha=P_{\alpha,h}$ are coercive. The enlarged block need
not yet be positive. Let $B_\alpha$ be its cross block and put
$\mathcal K=F_\alpha^{-1/2}B_\alpha A_\alpha^{-1/2}$.
For a candidate new shift $\beta>0$, define
\begin{align}
 H_A&=A_\alpha^{-1/2}(A_\beta-A_\alpha)A_\alpha^{-1/2},\\
 H_F&=F_\alpha^{-1/2}(F_\beta-F_\alpha)F_\alpha^{-1/2},\\
 H_B&=F_\alpha^{-1/2}(B_\beta-B_\alpha)A_\alpha^{-1/2}.
\end{align}
Congruence by the reference diagonal metrics makes the enlarged
form at $\beta$ equal to
\[
 \begin{pmatrix}
 I+H_A&(\mathcal K+H_B)^*\\
 \mathcal K+H_B&I+H_F
 \end{pmatrix}.
\]
It is coercive if and only if $I+H_F$ is coercive and
\begin{equation}\label{eq:si-joint}
 \boxed{I+H_A-(\mathcal K+H_B)^*(I+H_F)^{-1}
                     (\mathcal K+H_B)\succ0.}
\end{equation}
Here $\succ0$ means a strictly positive operator lower bound.
This is a second Schur complement; its signed changes retain
improvements in either diagonal storage or cross coupling.
A simple sufficient condition is
\[
 H_A\succeq aI,\quad H_F\succeq bI,\quad a,b>-1,\qquad
 \norm{\mathcal K+H_B}^2<(1+a)(1+b).
\]
Replacing each signed change by its absolute norm is more
conservative and may discard the improvement that makes the step
possible. All these statements concern operators, not only their
finite compressions.

\subsection{Scope near zero shift and at increasing depth}

The causal Cayley symbol has the core expansion
\begin{equation}
 \frac2\omega\frac{1-K_\omega(p)}{1+K_\omega(p)}
 =a_0(p)+\omega^2\left(\frac{a_0''(p)}6
                         -\frac{a_0(p)^3}{12}\right)+O(\omega^4),
\end{equation}
where primes mean $p$ derivatives. To see the coefficient, write
$\log K_\omega=-\omega a_0-\omega^3a_0''/6+O(\omega^5)$
and expand the resulting hyperbolic tangent. The expression is
even in $\omega$; its use after inversion is on appropriate smooth
cores, not an operator-norm comparison with the unbounded central
form. The extra cubic term also distinguishes accumulated storage
from the instantaneous generator.

Positive-shift comparisons in \eqref{eq:si-joint} avoid asserting
uniform convergence on the central form domain. They still require
rigorous estimates of the full blocks and their weak directions.
An iteration must additionally prove
$L_{j+1}=L_j+h_j$, $\sum_jh_j=\infty$ and
$\omega_j\downarrow0$. The algebra does not supply those estimates.
Discarding output through a projection, assuming monotonicity
in shift, or extrapolating local extension cannot replace them.

\subsection{Separate arithmetic anchor and append results}
\label{subsec:si-ema-progress}

A separate computer-assisted calculation in the critical-path
investigation establishes, in the original arithmetic normalization,
\begin{align}\label{eq:si-ema-anchor}
 Q_{0,1/2}&\succeq I/40,\\
 D_{10^{-3},1/2}&\succeq\delta I,\qquad
 \delta=\frac{24999}{500000000}=0.000049998.
 \notag
\end{align}
The first inequality is on the central form domain; the defect
bound holds for every $L^2$ input. The calculation uses a finite
part of the fixed arithmetic exponential-moving-average (EMA)
tower as a lower operator, exact endpoint response, and bounds
on the entire complementary input space. Outward rational interval
records passed at 40 and 60 digits \cite{EMAAnchor}.
This imports the stated certificate and its hypotheses, rather
than reproduce its numerical proof here. Specialist mathematical
review remains outstanding. It closes the earlier finite-input
pilot gap and is not a new depth horizon.

For $L=1/2$, $h=1/20$, $\omega=10^{-3}$, causal restriction transfers
the defect floor to the smaller window. Reflection supplies the
new-output orientation $F_{\mathrm{out}}=I-ZZ^*$, so both diagonal
defects in \eqref{eq:si-gluing} are bounded below by $\delta I$.
The complete append response, writing old distance from the join
as $\rho$, is
\[
 (Yf)(t)=\int_0^L k_\omega(t+\rho)f(L-\rho)\dd\rho,
 \qquad 0<t<h.
\]
It retains the full arithmetic history, both pole contributions
and the fixed local normalization.

The subsequent append calculation proves
\begin{equation}\label{eq:si-scalar-append-obstruction}
 55<\frac{\norm Y}{\delta}<72.
\end{equation}
Hence the scalar sufficient estimate
$\norm{\mathcal C}\leq\norm Y/\delta$ cannot prove
$\norm{\mathcal C}<1$; \eqref{eq:si-scalar-append-obstruction}
is not a lower bound above one for $\norm{\mathcal C}$.
Even after 32 cosine modes on both sides, the raw
complement/complement norm divided by $\delta$ exceeds 46
\cite{CumulativeAppend}. Complete-output finite-input diagnostics
instead give normalized quotients near $0.801830639$.
In exact arithmetic such finite quotients are lower, not upper,
bounds on the full coupling norm. These diagnostics alone give
no all-input upper bound below one.

The subsequent operator calculation retains directional defect
energy and controls both full complementary spaces, establishing
internally at the same $(L,h,\omega)$
\[
 \norm{F_{\mathrm{out}}^{-1/2}YE^{-1/2}}<0.951
 \qquad\cite{CumulativeCoupling}.
\]
It uses a forward/backward energy factorization and a Carleman
operator bound for the singular corner, preserving the complete
history and the output-defect orientation. Its proof and rational
records remain in the critical-path investigation; independent
specialist review is outstanding. It establishes neither an
all-depth continuation nor a physical realization. No instantaneous
operator lower bound is exponentiated through the nonnormal flow.
Auxiliary smoothing loss and averaging variance are not arithmetic
positivity; the field resolution in Section~\ref{sec:trace-hierarchy}
likewise makes no contribution to this arithmetic estimate.
